% !TEX root = ../main.tex
\chapter{Overview and Anatomy of a Dynamical System}\label{ch:overview}
\begin{paracol}{2}

\begin{sources}
\begin{tabularx}{\linewidth}{@{}l l L@{}}
\toprule
\textbf{Video} & \textbf{Length} & \textbf{Content}\\
\midrule
\seg{1}{1}{V1.1 Data-Driven Dynamical Systems Overview} & 21:43 & The general system $\dot{\vb x}=\vb f(\vb x,t,\vb u;\boldsymbol\beta)$; from first principles to data; challenges (nonlinearity, unknown dynamics, dimension, chaos, noise, multiscale, uncertainty); uses (prediction, design, control, interpretability); techniques (regression, neural networks, genetic programming); Koopman coordinates; physics in the learning.\\
\seg{1}{2}{V1.2 The Anatomy of a Dynamical System} & 17:53 & The same anatomy in more detail: state, dynamics, time, control, parameters, measurements; the challenges as optimization problems; interpretable and generalizable models.\\
\bottomrule
\end{tabularx}

\smallskip
\textbf{Transcripts} \file{materials/transcripts/C01-V01.txt}, \file{C01-V02.txt}.
\textbf{Book} Brunton \& Kutz, \emph{Data-Driven Science and Engineering}, Chapter~7 introduction and
Sec.~7.1 ``Overview, motivations, and challenges''. The second video (recorded about three years after
the first) repeats much of the first: the two are merged here, and each point is tagged with the
video where it is said best.
\end{sources}

\section*{Motivation}
These two short videos open the series. They contain no algorithm: they fix the language used in all
the following chapters (state, dynamics, control, parameters, measurements) and they set the agenda,
that is, which problems of dynamical systems are still open and why data and machine learning might
help with them. The thesis is simple. For three centuries dynamical systems were \emph{derived}: one
wrote down Newton's law, a Lagrangian or a Hamiltonian, and obtained the equations of motion. For many
systems of interest today (the brain, the climate, an epidemic, a financial market) there is no such
first-principles equation, but there is more and more \emph{data}. The program of the course is to
\emph{learn} the dynamical system from the data, in a form that is still interpretable
\ts{1}{1}{0:22}\ts{1}{1}{4:50}.

The reading for a physicist: nothing in this chapter is new mathematics, but the point of view is
not the textbook one. In a course on classical mechanics the equations are given and one studies
their solutions; here the solutions (the data) are given and one looks for the equations.

% ------------------------------------------------------------------
\section{The anatomy of a dynamical system}\label{sec:ov-anatomy}

\paragraph{Why data, why now.} Two things happened at the same time \ts{1}{1}{0:43}:
\begin{enumerate}
  \item an explosion of data, from experiments, simulations and historical records, in climate
        science, neuroscience, disease modelling, fluid dynamics and almost every other field;
  \item much better algorithms in machine learning, regression and optimization.
\end{enumerate}
Their confluence makes it possible, for the first time, to discover and characterize dynamical
systems purely from data \ts{1}{1}{1:03}.

\paragraph{The general form.} Dynamical systems are our models of the evolving world: how fluids mix,
how a rocket flies and lands, how a disease spreads on a network \ts{1}{2}{0:06}. They are
fundamentally tied to \emph{time}: they describe how things change and co-evolve in time
\ts{1}{2}{1:28}. The general object of the course is a system of coupled ordinary differential
equations \ts{1}{1}{2:26}\ts{1}{2}{1:49}
\begin{equation}\label{eq:ov-general}
  \frac{\mathrm d}{\mathrm dt}\vb x(t) = \vb f\big(\vb x(t),\,t,\,\vb u(t);\,\boldsymbol\beta\big),
  \qquad
  \vb y(t) = \vb g\big(\vb x(t),\,t\big),
\end{equation}
together with the measurements $\vb y$ that we can actually take of it. Each symbol is one ``organ''
of the anatomy.

\begin{definition}[New concepts: the anatomy of a dynamical system]\label{def:ov-anatomy}
In \eqref{eq:ov-general}:
\begin{enumerate}[(a)]
  \item the \textbf{state} $\vb x\in\R^n$ is the minimal set of numbers that describes the system
        uniquely at one instant; $n$ is the \emph{dimension} of the system \ts{1}{2}{1:49};
  \item the \textbf{dynamics} $\vb f=(f_1,\dots,f_n)$ is a vector-valued function: $f_j$ gives the
        rate of change $\dot x_j$ of the $j$-th state variable. Since it assigns a vector
        $\vb f(\vb x)$ to every point $\vb x$ of state space, it is also called a \textbf{vector field}
        \ts{1}{1}{3:49}\ts{1}{2}{2:52};
  \item $t$ is \textbf{time}. If $\vb f$ depends explicitly on $t$ the system is
        \textbf{non-autonomous}, otherwise it is \textbf{autonomous} \ts{1}{1}{2:26};
  \item the \textbf{control} (or \emph{actuation}, \emph{input}) $\vb u$ collects the variables we can
        manipulate actively, the ``knobs'' \ts{1}{2}{4:14};
  \item the \textbf{parameters} $\boldsymbol\beta$ are all the other constants on which $\vb f$
        depends, which we do not control but want to understand \ts{1}{2}{4:55};
  \item the \textbf{measurements} (or \emph{outputs}) $\vb y=\vb g(\vb x,t)$ are what we can actually
        observe; in general $\vb y$ has far fewer components than $\vb x$ \ts{1}{2}{5:15}.
\end{enumerate}
\end{definition}

\begin{aside}[Where the names come from]
\emph{Dynamical} comes from the Greek \emph{dýnamis}, ``power, force'': dynamics was the part of
mechanics that studies motion together with the forces that cause it (Leibniz coined
\emph{dynamica} in the 1690s). The word survived the forces: a ``dynamical system'' today is any rule
that makes a state evolve in time. \emph{Autonomous} is from the Greek \emph{autós} ``self'' and
\emph{nómos} ``law'': the system obeys its own law, with no external clock.
\emph{Actuation} is from the medieval Latin \emph{actuare}, ``to put into action''.
\end{aside}

\paragraph{The state.} The state is the minimal description: for a pendulum the angle $\theta$ alone
is not enough, because the same angle can be crossed with different speeds; the state is
$\vb x=(\theta,\dot\theta)$ \ts{1}{2}{2:09}. The state can be enormous: for the stock market it might be
the value of all the stocks (and possibly more information), for the weather the conditions on a
one-kilometre grid covering the Earth \ts{1}{2}{2:31}.

\begin{example}[The pendulum]\label{ex:ov-pendulum}
A pendulum of length $L$ in gravity $g$ obeys $\ddot\theta=-(g/L)\sin\theta$. With
$\vb x=(x_1,x_2)=(\theta,\dot\theta)$ it becomes a first-order system
\[
  \dot x_1=x_2,\qquad \dot x_2=-\frac{g}{L}\sin x_1 ,
\]
so $n=2$ and $\boldsymbol\beta=(g,L)$ (and the mass, which drops out here). If the pivot is mounted on
a cart that we can drive, the cart's acceleration is a control $u$ \ts{1}{2}{4:34}. If we film the
pendulum and extract only the horizontal position of the bob, the measurement is
$y=L\sin x_1$: one number, from which the state must be reconstructed.
\end{example}

\paragraph{The dynamics as a vector field.} Every $n$-th order ODE can be written as a first-order
system by enlarging the state, as in the pendulum; so \eqref{eq:ov-general} is fully general. The
picture to keep in mind is the vector field: at every point of state space $\vb f$ is an arrow
saying how fast and in which direction the state moves \ts{1}{1}{4:09}. A solution is a curve
everywhere tangent to the arrows (\cref{fig:ov-pendulum}). Often $\vb f$ comes from Newton's second
law or something like it \ts{1}{2}{3:12}.

\begin{nfigure}
\begin{tikzpicture}[x=1.15cm,y=0.95cm,>=Stealth]
  % vector field of the pendulum, arrows normalised to a fixed length
  \foreach \t in {-3.0,-2.5,...,3.01} {
    \foreach \w in {-2.5,-2.0,...,2.51} {
      \pgfmathsetmacro{\u}{\w}
      \pgfmathsetmacro{\v}{-sin(\t r)}
      \pgfmathsetmacro{\n}{max(sqrt(\u*\u+\v*\v),0.001)}
      \draw[->,srccol!70,thin] (\t,\w) -- ++(0.32*\u/\n,0.32*\v/\n);
    }
  }
  % closed orbits (oscillations): energy E = w^2/2 - cos(theta) = c
  \foreach \c in {-0.8,-0.4,0.3} {
    \pgfmathsetmacro{\tm}{acos(-\c)*pi/180-0.005}
    \draw[intcol,thick,smooth,samples=60,domain=-\tm:\tm] plot (\x,{sqrt(2*(\c+cos(\x r)))});
    \draw[intcol,thick,smooth,samples=60,domain=-\tm:\tm] plot (\x,{-sqrt(2*(\c+cos(\x r)))});
  }
  % separatrix E = 1 and two rotations E = 1.6
  \draw[thmcol,thick,smooth,samples=80,domain=-pi:pi] plot (\x,{2*cos(\x/2 r)});
  \draw[thmcol,thick,smooth,samples=80,domain=-pi:pi] plot (\x,{-2*cos(\x/2 r)});
  \draw[defcol,thick,smooth,samples=80,domain=-pi:pi] plot (\x,{sqrt(2*(1.6+cos(\x r)))});
  \draw[defcol,thick,smooth,samples=80,domain=-pi:pi] plot (\x,{-sqrt(2*(1.6+cos(\x r)))});
  % axes and fixed points
  \draw[->] (-3.4,0) -- (3.6,0) node[right] {$\theta$};
  \draw[->] (0,-2.9) -- (0,3.1) node[above] {$\dot\theta$};
  \foreach \t/\l in {-pi/{-\pi},pi/{\pi}} \draw (\t,0.08) -- (\t,-0.08) node[below] {$\l$};
  \fill[intcol] (0,0) circle (2pt);
  \fill[thmcol] (-pi,0) circle (2pt) (pi,0) circle (2pt);
\end{tikzpicture}
\caption{State space of the pendulum $\dot\theta=\omega$, $\dot\omega=-\sin\theta$ ($g/L=1$). Grey: the
vector field $\vb f$ (arrows of equal length, only the direction is shown). Green: oscillations
around the stable equilibrium; red: the separatrix through the unstable (upright) equilibria; blue:
rotations, the pendulum goes over the top. Each curve is a level set of the energy
$E=\tfrac12\omega^2-\cos\theta$.}\label{fig:ov-pendulum}
\end{nfigure}

\paragraph{Time.} The explicit dependence on $t$ matters, because in many systems the rules themselves
change in time \ts{1}{2}{3:33}. The equations that describe your brain while you watch a lecture are
different from those that describe it while you sleep; the weather and the climate change; we may be
actively changing the dynamics of a system, as with the ocean \ts{1}{2}{3:53}.

\paragraph{Control versus parameters.} Both $\vb u$ and $\boldsymbol\beta$ are ``inputs'' of $\vb f$,
but they play opposite roles. The control $\vb u$ is what we act on, possibly in real time; the
parameters $\boldsymbol\beta$ are given by the situation. For an aircraft, $\boldsymbol\beta$ would
contain wind speed, humidity, altitude and weight, and $\vb u$ the deflection of the flaps and of the
elevators \ts{1}{1}{3:28}. We want to understand the explicit dependence of $\vb f$ on
$\boldsymbol\beta$, because a small change of a parameter can produce a large qualitative change in the
dynamics, a \emph{bifurcation} (here it is only named; bifurcations are treated in Chapter~2) \ts{1}{2}{4:55}.

\paragraph{Measurements.} We rarely have access to the full state \ts{1}{2}{5:15}. One cannot record
every neuron of a brain; at best a small subset, or a crude proxy such as an EEG or an ECG
\ts{1}{2}{5:36}. Hence the second equation in \eqref{eq:ov-general}. The data from which models are
learned are samples of $\vb y$, $t$, $\vb u$ and possibly $\boldsymbol\beta$ \ts{1}{1}{4:29}.

\begin{nfigure}
\begin{tikzpicture}[>=Stealth,
    blk/.style={draw=defcol,very thick,fill=defcol!7,rounded corners=2pt,minimum width=4.6cm,
      minimum height=1.25cm,align=center},
    lab/.style={font=\small,align=center}]
  \node[blk] (sys) {$\dot{\vb x}=\vb f(\vb x,t,\vb u;\boldsymbol\beta)$\\[2pt]\small state $\vb x\in\R^n$, possibly $n\sim10^9$};
  \node[blk,draw=intcol,fill=intcol!7,minimum width=2.6cm,right=1.2cm of sys] (meas) {$\vb y=\vb g(\vb x,t)$};
  \draw[->,very thick] ([xshift=-1.4cm]sys.west) node[left,lab] {control $\vb u$\\(knobs\\we turn)} -- (sys.west);
  \draw[->,very thick] ([yshift=1.2cm]sys.north) -- node[right,lab] {parameters $\boldsymbol\beta$ (given; bifurcations)} (sys.north);
  \draw[->,very thick] (sys.east) -- node[above,lab] {$\vb x(t)$} (meas.west);
  \draw[->,very thick] (meas.east) -- ++(1.0,0) node[right,lab] {data\\(few,\\noisy)};
  \draw[->,thick,dashed,srccol] ([yshift=-0.9cm]sys.south) -- node[right,lab,text=srccol] {disturbances, noise} (sys.south);
\end{tikzpicture}
\caption{The anatomy of a dynamical system as a block diagram, the way control engineers draw it
(Chapter~8 of the book). The machine-learning tasks of the course attack one block at a time: learn
$\vb f$, reconstruct $\vb x$ from $\vb y$, choose $\vb u$, study the dependence on
$\boldsymbol\beta$.}\label{fig:ov-anatomy}
\end{nfigure}

\begin{secondary}[Precise setting: manifolds, existence and uniqueness]
In the book the autonomous system $\dot{\vb x}=\vb f(\vb x)$ is formulated more generally: the state
lives on a smooth $n$-dimensional manifold $M$ and $\vb f$ is a vector field, a section of the tangent
bundle $TM$, so that $\vb f(\vb x)\in T_{\vb x}M$. Example: the pendulum angle is defined modulo
$2\pi$, so the true state space is the cylinder $S^1\times\R$, and \cref{fig:ov-pendulum} should be
rolled up so that $\theta=-\pi$ and $\theta=\pi$ coincide. In practice one takes $M=\R^n$ and asks
$\vb f$ to be \emph{Lipschitz}, which guarantees existence and uniqueness of solutions.
\begin{theorem}[Picard--Lindel\"of]\label{thm:ov-picard}
Let $\vb f:\R\times\R^n\to\R^n$ be continuous in $t$ and Lipschitz in $\vb x$ on a neighbourhood of
$(t_0,\vb x_0)$: $\lVert\vb f(t,\vb x)-\vb f(t,\vb x')\rVert\le L\lVert\vb x-\vb x'\rVert$. Then the
problem $\dot{\vb x}=\vb f(t,\vb x)$, $\vb x(t_0)=\vb x_0$ has a unique solution on some interval
$[t_0-\epsilon,t_0+\epsilon]$.
\end{theorem}
Uniqueness is what makes ``state'' a meaningful word: two trajectories cannot cross, so the present
state determines the future. It fails for $\dot x=\sqrt{|x|}$ at $x=0$ (not Lipschitz): both $x(t)=0$
and $x(t)=t^2/4$ start from $x(0)=0$.
\end{secondary}

\begin{exercise}[Non-autonomous to autonomous]
Show that $\dot{\vb x}=\vb f(\vb x,t)$, $\vb x\in\R^n$, can be written as an autonomous system in
$\R^{n+1}$. Apply this to the forced oscillator $\ddot x+\gamma\dot x+\omega_0^2x=F\cos\Omega t$, and
then show that, because the forcing is periodic, the extra coordinate can be taken on a circle
$S^1$ instead of $\R$. What is the dimension of the state space of the forced pendulum?
\end{exercise}

\begin{exercise}[Eigenvalues of the damped oscillator]
Write $\ddot x+2\zeta\omega_0\dot x+\omega_0^2x=0$ as $\dot{\vb x}=A\vb x$.
\begin{enumerate}[(a)]
  \item Compute the eigenvalues of $A$ and classify the motion for $0<\zeta<1$, $\zeta=1$, $\zeta>1$.
  \item For $\zeta=1$ the matrix is not diagonalizable: compute $e^{At}$ anyway (Jordan form) and
        find the $te^{-\omega_0t}$ term.
  \item Sample the underdamped solution with step $\Delta t$: give the eigenvalues $\mu_\pm$ of the
        discrete-time map and show that $|\mu_\pm|<1$. Where do they lie in the complex plane as
        $\Delta t$ grows?
\end{enumerate}
\end{exercise}

% ------------------------------------------------------------------
\section{From first principles to data}\label{sec:ov-data}

\paragraph{The classical way.} Dynamical systems are not new: Newton's second law $F=ma$ is already a
dynamical system, and it gives the trajectories of projectiles, the motion of a pendulum and of the
planets \ts{1}{1}{4:50}. For centuries the method has been \emph{physics based}: write down Newton's
law, the Lagrangian or the Hamiltonian, and derive the equations of motion from first principles
\ts{1}{1}{5:11}\ts{1}{2}{6:17}. With these models we built satellites and went to the Moon
\ts{1}{1}{6:53}.

\begin{history}[Newton, Poincar\'e and the birth of dynamical systems]
The book dates \emph{modern} dynamical systems to Henri Poincar\'e's work on the motion of planets.
In 1889 Poincar\'e won the prize offered by King Oscar~II of Sweden for a memoir on the stability of
the solar system (the three-body problem). While the memoir was being printed, queries from the editor
Lars Phragm\'en led Poincar\'e to a serious error; correcting it, he discovered that the stable and unstable manifolds of a periodic
orbit can intersect transversally (homoclinic tangle), the first glimpse of chaos. The printed copies
were recalled and destroyed, and Poincar\'e paid for the reprinting, more than the 2500 crowns of the
prize. The corrected memoir appeared in \emph{Acta Mathematica} in 1890 and grew into \emph{Les
m\'ethodes nouvelles de la m\'ecanique c\'eleste} (1892--1899), where the qualitative, geometric
study of ODEs begins.\par
\textit{References:} J.~Barrow-Green, \emph{Poincar\'e and the Three Body Problem} (AMS, 1997);
Brunton \& Kutz, Ch.~7 introduction.
\end{history}

\paragraph{Why it is no longer enough.} For the systems we most want to understand today, such as the
brain, the climate, a financial market or the spread of a disease, there are no first-principles
equations that are easy to write, simulate and control \ts{1}{1}{5:52}. ``There is no master equation
for your brain'', and no agreed governing law for how a disease spreads or how the climate works
\ts{1}{1}{6:13}\ts{1}{2}{6:57}. But there are more and more data: recordings of brain activity,
for example, come in ever larger volumes \ts{1}{1}{6:13}. The next generation of problems, like
controlling turbulence or understanding the brain, will not always be solvable with first-principles
models, so the models will have to be derived from data \ts{1}{1}{7:14}. The field is in its infancy:
knowing the physics still lets you do much more than data alone \ts{1}{1}{7:55}.

\begin{remark}[All models are models]
Even $F=ma$, the Navier--Stokes equations of fluids and the Schr\"odinger equation are not ``the
truth'' but \emph{models} of reality, excellent ones \ts{1}{1}{6:33}. Seen this way, learning a model
from data is not a different activity from physics but the same one, done with other tools. The book
adds that, as systems become more complex, choosing the ``correct'' model becomes more subjective,
which is a reason to want automatic, reproducible model discovery.
\end{remark}

\paragraph{Reduced-order models.} Data-driven models are useful even where the equations are known.
The Navier--Stokes equations describe a fluid exactly but may be far too expensive to use; what one
wants is a \emph{reduced-order model}, a dynamical system for a much smaller state that captures only
what we care about in the flow \ts{1}{2}{7:18}.

\begin{definition}[New concepts: data-driven and reduced-order models]\label{def:ov-models}
Two ways of obtaining a model, and one kind of model:
\begin{enumerate}[(a)]
  \item A \textbf{first-principles} model is derived from accepted physical laws (Newton, Lagrange,
        Hamilton, conservation laws); a \textbf{data-driven} model is obtained by fitting
        (regression, optimization) to measurements of the system.
  \item A \textbf{reduced-order model} (ROM) is a dynamical system $\dot{\vb a}=\tilde{\vb f}(\vb a)$
        for a state $\vb a\in\R^r$ with $r\ll n$, built so that a known map $\vb a\mapsto\vb x$
        reproduces the quantities of interest of the full system $\dot{\vb x}=\vb f(\vb x)$.
\end{enumerate}
\end{definition}

% ------------------------------------------------------------------
\section{Challenges of modern dynamical systems}\label{sec:ov-challenges}

Brunton's list is explicitly personal: another dynamicist would give a different one, and this is his
group's ``to-do list'' \ts{1}{1}{12:03}. The two primary challenges, those that data can help with most,
are \ts{1}{1}{9:38}: nonlinearity and unknown dynamics.

\paragraph{1. Nonlinearity.} If the system is linear, $\dot{\vb x}=A\vb x$, everything is known
from the eigenvalues and eigenvectors of $A$ \ts{1}{1}{8:16}. Even a small nonlinearity, an $x^2$, an
$x^3$ or a $\cos x$ on the right-hand side, makes it very hard to understand, predict, estimate and
control the future \ts{1}{1}{8:57}. Superposition fails: the sum of two solutions is no longer a
solution \ts{1}{2}{7:38}. For linear systems ``you could just write down the solution and go home,
and we would all be out of business'': nonlinearity is ``job security'' \ts{1}{2}{7:59}. There are
special classes of nonlinear systems and special solutions, but no general, one-size-fits-all theory
\ts{1}{2}{8:19}. Why linear systems count as ``solved'' is worth seeing in detail.

\begin{proposition}[Linear systems are solved by their eigenvectors]\label{prop:ov-linear}
Let $\dot{\vb x}=A\vb x$ with $A\in\R^{n\times n}$ diagonalizable, $AT=T\Lambda$, where
$\Lambda=\diag(\lambda_1,\dots,\lambda_n)$ and the columns $\boldsymbol\xi_j$ of $T$ are eigenvectors.
Then
\begin{enumerate}[(i)]
  \item the solution is $\vb x(t)=e^{At}\vb x(0)=T e^{\Lambda t}T^{-1}\vb x(0)
        =\sum_{j=1}^n b_j e^{\lambda_j t}\boldsymbol\xi_j$, with $\vb b=T^{-1}\vb x(0)$;
  \item in the eigenvector coordinates $\vb z=T^{-1}\vb x$ the dynamics decouple: $\dot z_j=\lambda_j z_j$;
  \item sampled at times $t_k=k\Delta t$, the system is the linear map $\vb x_{k+1}=e^{A\Delta t}\vb x_k$,
        whose eigenvalues are $\mu_j=e^{\lambda_j\Delta t}$.
\end{enumerate}
\end{proposition}
\begin{proof}
Substitute $\vb x=T\vb z$: $T\dot{\vb z}=AT\vb z=T\Lambda\vb z$, and multiplying by $T^{-1}$ gives
$\dot{\vb z}=\Lambda\vb z$, i.e.\ (ii). Each scalar equation has solution $z_j(t)=e^{\lambda_jt}z_j(0)$,
so $\vb z(t)=e^{\Lambda t}\vb z(0)$ and $\vb x(t)=Te^{\Lambda t}T^{-1}\vb x(0)$; writing $T e^{\Lambda t}\vb b$
column by column gives the sum in (i). The power series of the exponential gives
$e^{At}=\sum_k (T\Lambda T^{-1})^kt^k/k!=T(\sum_k\Lambda^kt^k/k!)T^{-1}=Te^{\Lambda t}T^{-1}$.
For (iii), $\vb x_{k+1}=e^{A\Delta t}\vb x_k$ follows from (i) with initial time $t_k$, and
$e^{A\Delta t}T=Te^{\Lambda\Delta t}$ shows that the eigenvectors are the same, with eigenvalues
$e^{\lambda_j\Delta t}$.
\end{proof}

\begin{intuition}[Normal modes]
\Cref{prop:ov-linear} is the normal-mode analysis of small oscillations in classical mechanics. Each
eigenvector $\boldsymbol\xi_j$ is a spatial pattern that evolves on its own, growing or decaying at rate
$\operatorname{Re}\lambda_j$ and oscillating at frequency $\operatorname{Im}\lambda_j$; any motion is
a superposition of them. The whole course can be read as the search for such modes when $\vb f$ is
nonlinear or unknown: DMD (Chapter~4) fits them to data, and Koopman theory (Chapter~6) looks for
coordinates in which a nonlinear system has them.
\end{intuition}

\paragraph{2. Unknown dynamics.} Perhaps more fundamental: often $\vb f$ is simply not known
\ts{1}{1}{9:18}. This challenge ``goes unmentioned in a lot of cases'' \ts{1}{2}{6:17}. For a pendulum
you open a physics book and derive the equations; for the brain or a disease network we have models
and approximations, but no agreed governing law, and $\vb f$ must be discovered from data
\ts{1}{2}{6:37}\ts{1}{2}{6:57}. Even when $\vb f$ is known it may be unusable in practice, as with the
Navier--Stokes equations, and one needs a reduced model \ts{1}{2}{7:18}.

\paragraph{Other challenges.} Important, but secondary to the first two \ts{1}{1}{9:58}:
\begin{itemize}
  \item \textbf{High dimension.} The state may have millions or billions of components. To simulate
        the climate one chops the Earth into maybe a million to a billion grid cells, at a resolution
        of a kilometre or less \ts{1}{1}{10:20}; a brain has about a hundred billion neurons, so a
        hundred-billion-dimensional state just to say which neurons are active \ts{1}{2}{8:40}; a
        realistic epidemic model needs the network of contacts of billions of people
        \ts{1}{2}{9:00}.
  \item \textbf{Multiscale dynamics.} Coherent structures at many scales of space and time
        \ts{1}{1}{11:21}. In a turbulent fluid you can zoom in, and in, and still find structure
        \ts{1}{1}{11:42}. For the weather, Earth-scale seasonal patterns and hurricanes matter, but so
        do tiny features, for a forecast three weeks out \ts{1}{2}{9:41}. The brain is multiscale in
        space (regions made of neurons) and in time: you hear words over seconds, connect them to what
        was said minutes ago, against a lifetime of experience \ts{1}{2}{10:02}.
  \item \textbf{Chaos and transients.} Not all nonlinear systems are chaotic, but many are: the
        future depends sensitively on small changes of the initial condition, of the parameters and of
        the dynamics \ts{1}{2}{10:22}. This is why the weather cannot be predicted three weeks ahead:
        the system is chaotic and we can measure it, and know its equations and parameters, only so
        accurately \ts{1}{2}{11:02}.
  \item \textbf{Latent (hidden) variables.} With a big state, we are surely not measuring everything:
        a few neurons, or a proxy. The unmeasured variables may matter a lot \ts{1}{2}{11:23}.
  \item \textbf{Noise and stochasticity.} Real systems are forced by disturbances and measured with
        noise \ts{1}{1}{11:01}\ts{1}{2}{11:45}. And the noise is not the Gaussian white noise of
        textbooks (``the only thing we can actually solve'') but correlated, structured and biased
        \ts{1}{2}{11:45}.
  \item \textbf{Uncertainty} in everything: the initial condition, the parameters, the model $\vb f$
        itself, what is measured and how, even \emph{when} a measurement was taken; it must be handled
        and propagated \ts{1}{1}{12:03}\ts{1}{2}{12:05}.
\end{itemize}
The systems of interest, such as neuroscience, turbulence, financial markets and disease, have all of
these at once: high-dimensional, multiscale, nonlinear, with unknown $\vb f$ \ts{1}{1}{12:24}.

\begin{nfigure}
\begin{tikzpicture}[>=Stealth,
    main/.style={draw=thmcol,very thick,fill=thmcol!8,rounded corners=2pt,minimum width=3.2cm,
      minimum height=0.9cm,align=center,font=\small\bfseries},
    sec/.style={draw=srccol,fill=srccol!8,rounded corners=2pt,minimum width=1.95cm,minimum height=0.85cm,
      align=center,font=\footnotesize},
    use/.style={draw=intcol,thick,fill=intcol!8,rounded corners=2pt,minimum width=3.0cm,
      minimum height=0.75cm,align=center,font=\footnotesize}]
  \node[main] (nl) at (-2.2,0) {nonlinearity};
  \node[main] (unk) at (2.2,0) {unknown $\vb f$};
  \node[draw,very thick,keycol,circle,minimum size=1.4cm,align=center,font=\footnotesize\bfseries,text=black]
    (opt) at (0,-1.7) {optim-\\ization};
  \node[sec] (ms) at (-4.4,-3.4) {multiscale};
  \node[sec] (dim) at (-2.2,-3.4) {high\\dimension};
  \node[sec] (ch) at (0,-3.4) {chaos,\\transients};
  \node[sec] (lat) at (2.2,-3.4) {latent\\variables};
  \node[sec] (noise) at (4.4,-3.4) {noise,\\uncertainty};
  \foreach \n in {nl,unk,ms,dim,ch,lat,noise} \draw[->,srccol] (\n) -- (opt);
  \node[use] (u1) at (7.6,0.2) {prediction\\(ensembles)};
  \node[use,below=0.3cm of u1] (u2) {design and\\optimization};
  \node[use,below=0.3cm of u2] (u3) {estimation and\\control};
  \node[use,below=0.3cm of u3] (u4) {interpretability,\\generalizability};
  \draw[->,very thick,keycol] (opt.east) to[out=0,in=180] node[below,font=\footnotesize,text=black,pos=0.4] {models} (u3.west);
  \node[font=\small\itshape] at (0,0.85) {challenges};
  \node[above=0.1cm of u1,font=\small\itshape] {uses};
\end{tikzpicture}
\caption{Map of the chapter. The two primary challenges (red) and the secondary ones (grey) can all be
recast as optimization problems, which is where machine learning enters \ts{1}{2}{12:47}; the
resulting models serve the four uses (green).}\label{fig:ov-map}
\end{nfigure}

\paragraph{Challenges as optimization problems.} Most of these challenges can be recast as
\emph{optimization} problems \ts{1}{2}{12:47}: handle noise and disturbances optimally in a prediction
or estimate, propagate uncertainty optimally, reconstruct the full state optimally from the limited
measurements $\vb y$. And machine learning is, at bottom, ``optimization based on data'': building
models by optimization over a wealth of data \ts{1}{2}{13:07}. So the anatomy of
\cref{def:ov-anatomy} lets one pick a single task and attack it with these tools \ts{1}{2}{13:28}.
Every challenge on the list is not one problem but whole careers of work \ts{1}{2}{14:29}.

\begin{example}[State reconstruction as optimization]
Given measurements $\vb y=C\vb x+\boldsymbol\eta$ with $C\in\R^{p\times n}$, $p\ll n$, the state cannot be
obtained by inverting $C$. One instead solves an optimization problem: least squares with a prior
(the Kalman filter, Chapter~8 of the book), or least squares in a basis where $\vb x$ is sparse
(compressed sensing, Chapter~4 of these notes).
\end{example}

\begin{exercise}[How big is the weather?]
Estimate the dimension $n$ of the state of a global weather model with a horizontal resolution of
$1\,$km, $100$ vertical levels and $6$ variables per cell (three velocity components, temperature,
pressure, humidity); the Earth's surface is about $5.1\times10^8\,\mathrm{km}^2$. How many bytes does
one snapshot take in double precision? How many snapshots of this size fit in a petabyte?
\end{exercise}

\begin{exercise}[Superposition]
Show that the set of solutions of $\dot{\vb x}=A\vb x$ is a vector space. Then show, with an explicit
counterexample, that for $\dot x=-x+x^2$ the sum of two solutions is not a solution. Solve this
equation exactly (separation of variables or the substitution $\varphi=1/x$) and discuss which initial
conditions blow up in finite time.
\end{exercise}

% ------------------------------------------------------------------
\section{What models are for}\label{sec:ov-uses}

Again a personal, not definitive, list \ts{1}{1}{13:06}.

\paragraph{1. Future state prediction.} The most obvious use \ts{1}{1}{13:06}\ts{1}{2}{13:48}:
forecast the weather a week ahead, or, if someone is lost at sea, integrate from a rough initial
position where they will be two hours later. Prediction need not be of a single trajectory: it can be
statistical, an \emph{ensemble} of initial conditions propagated together to give a probability
density of the future state \ts{1}{1}{13:26}\ts{1}{2}{14:09}. For the climate we care not about one
trajectory but about the probability of evolving towards higher or lower temperatures
\ts{1}{1}{13:47}; for the Gulf oil spill, the source is known only up to a Gaussian, which must be
propagated into a forecast of where the oil will be \ts{1}{2}{14:09}.

\paragraph{2. Design and optimization.} With a model, tune the parameters $\boldsymbol\beta$ to get
a desired performance \ts{1}{1}{14:09}\ts{1}{2}{14:51}. Formula One cars, America's Cup yachts and
rocket engines are complex multiphysics systems dominated by turbulent fluid dynamics; the dynamics
change a lot with rain, humidity or the shape of an aerodynamic scoop, and so does the drag. Much
effort goes into reduced-order models of these systems, cheap enough to optimize the design rapidly
\ts{1}{1}{14:29}\ts{1}{1}{14:51}. This requires understanding how the system depends on its parameters
\ts{1}{2}{15:11}.

\paragraph{3. Control.} Different from design: here we measure the system in real time and, with a
model, compute an actuation signal that modifies its behaviour through sensor-based
\emph{feedback} \ts{1}{1}{15:11}. If design optimization is done in real time from measurements,
the parameters become actuators and design becomes feedback control \ts{1}{2}{15:11}. Examples:
suppressing epileptic seizures, suppressing a disease as it spreads through a continent, increasing the
power extracted by a turbine from a flow \ts{1}{1}{15:32}.

\paragraph{4. Interpretability and physical intuition.} Often missing from such lists, but for
Brunton maybe the most important \ts{1}{1}{15:52}. Many machine-learning methods are hard to interpret
and give no intuition; Brunton argues for tailoring them so that the models they produce are
interpretable and teach us more physics than we knew \ts{1}{1}{16:15}. Besides manipulating and
predicting the world, we want to understand how and why it works \ts{1}{2}{15:53}. Two key words
\ts{1}{2}{16:13}:
\begin{itemize}
  \item \textbf{interpretable}: the model is simple enough that we can communicate and reason about it;
  \item \textbf{generalizable}: the model works in situations it was not built for, for instance at
        parameter values never tested \ts{1}{2}{16:33}.
\end{itemize}
Newton's second law is the ultimate interpretable and generalizable model: so simple that anyone can
understand it, valid in a vast range of situations, and it changed how we understand and interact with
the world \ts{1}{2}{16:33}\ts{1}{2}{16:54}.

\begin{intuition}[Interpretable = few terms]
For a physicist, ``interpretable'' has a concrete meaning that will come back in SINDy (Chapter~5): a
model with few terms, each one recognizable (a damping, a quadratic coupling, a restoring force). A
neural network with $10^6$ weights may predict well, but it does not tell you that $\vb f$ contains
$-\sin\theta$. Generalization is the reward for parsimony: a law with few terms extrapolates, a
curve fitted with many terms does not (overfitting).
\end{intuition}

% ------------------------------------------------------------------
\section{Techniques and the plan of the course}\label{sec:ov-techniques}

\paragraph{Methods.} An incomplete list of the techniques that work \ts{1}{1}{16:57}:
\begin{itemize}
  \item \textbf{regression}, linear and \emph{sparse}, probably the foremost one \ts{1}{1}{17:17};
  \item \textbf{neural networks} and deep learning, very powerful at learning dynamics, but with
        interpretability problems; the course shows how to constrain them so that they are interpretable
        by construction \ts{1}{1}{17:39};
  \item \textbf{genetic algorithms} and \textbf{genetic programming}, used effectively to discover
        dynamical systems from data \ts{1}{1}{17:39}.
\end{itemize}

\begin{definition}[New concepts: sparse regression, genetic programming]
Two ways of finding a model with few terms:
\begin{enumerate}[(a)]
  \item \textbf{Sparse regression} fits $\vb y\approx\Theta\boldsymbol\xi$ while asking that most
        entries of the coefficient vector $\boldsymbol\xi$ vanish (by an $\ell_1$ penalty or by
        thresholding), so that only a few columns of $\Theta$, the few ``terms'' of the model,
        survive. \emph{Sparse}, from the Latin \emph{sparsus}, ``scattered'': few non-zero entries
        scattered in a sea of zeros.
  \item \textbf{Genetic programming} searches over formulas (expression trees) by mimicking evolution:
        a population of candidate equations is mutated and recombined, and the ones that fit the data
        best with the fewest terms survive.
\end{enumerate}
\end{definition}

\begin{aside}[Symbolic regression by evolution]
The best-known example is the work of M.~Schmidt and H.~Lipson, \emph{Distilling free-form natural
laws from experimental data}, Science 324 (2009), who recovered Hamiltonians and Lagrangians of
double pendulums and oscillators from motion-tracking data with genetic programming (ref.~[477] of
the book's Chapter~7, together with Bongard and Lipson 2007, ref.~[68]).
\end{aside}

\paragraph{Coping with nonlinearity: Koopman.} A preview of a large part of the course
\ts{1}{1}{18:21}. Koopman analysis looks for a nonlinear change of coordinates $\boldsymbol\varphi(\vb x)$
such that, in the new coordinates, the dynamics is linear,
\[
  \frac{\mathrm d}{\mathrm dt}\boldsymbol\varphi(\vb x) = K\,\boldsymbol\varphi(\vb x),
\]
so that the whole linear toolbox of \cref{prop:ov-linear} applies. Finding coordinates that make a
system simpler is an old goal of physics; what is new is using deep learning or sparse regression to
find them from data \ts{1}{1}{18:41}\ts{1}{1}{19:02}. The other half of the course is to discover $\vb f$
itself by regression \ts{1}{1}{19:02}. The methods to come: dynamic mode decomposition, Koopman
analysis, delay coordinates, neural networks \ts{1}{1}{19:23}.

\begin{example}[A nonlinear system made linear]
$\dot x=x^2$ is nonlinear. With $\varphi(x)=1/x$ (for $x\neq0$), $\dot\varphi=-\dot x/x^2=-1$: in
the coordinate $\varphi$ the motion is uniform, $\varphi(t)=\varphi(0)-t$, so
$x(t)=x_0/(1-x_0t)$, with the blow-up at $t=1/x_0$ in plain view. Chapter~7 shows that such
finite linear representations exist only in special cases.
\end{example}

\paragraph{Put the physics in.} The final message \ts{1}{1}{19:43}. Most of machine learning was
developed for \emph{static} tasks, like image classification or scene labelling. Dynamical systems
instead have physics: constraints, conservation of energy and mass, $F=ma$
\ts{1}{1}{20:03}. Even where the exact equations are unknown, as for a reduced model of a turbulent
fluid, the brain or an epidemic, there are surely constraints, symmetries and invariants
\ts{1}{1}{20:03}\ts{1}{1}{20:25}. Building them into the learning procedure pays twice: the models are
more interpretable, and they can be learned with less data \ts{1}{1}{20:45}\ts{1}{1}{21:06}.

\begin{intuition}[Symmetries as priors]
Enforcing a conservation law shrinks the space of candidate models. If energy is conserved, $\vb f$
must satisfy $\nabla E\cdot\vb f=0$: a linear constraint on the coefficients of a regression, which
removes whole directions of the search space. Fewer free parameters, less data needed: this is the
statistical face of Noether's theorem.
\end{intuition}

\paragraph{Homework from the lecture.} Brunton closes with an invitation \ts{1}{2}{17:15}: pick a system
you care about (the brain, an epidemic, the climate) and ask what its state vector is, what its
controls and parameters are, what rules might govern its evolution, what can be measured, and where it
sits among the challenges and the uses \ts{1}{2}{17:35}. The exercise at the end of the chapter does exactly this.

\begin{exercise}[Energy constraint as a linear constraint]
Suppose you look for a model of a one-degree-of-freedom conservative system, $\dot q=p$,
$\dot p=\sum_{k=0}^{3}\xi_kq^k$, as a linear combination of monomials. Show that energy conservation is
automatic for any $\boldsymbol\xi$, and write the energy. Now allow also $\dot p$ to contain a term
$\xi_4p$: which value of $\xi_4$ does energy conservation impose? Relate this to the last intuition box
of \cref{sec:ov-techniques}.
\end{exercise}

\begin{keyformulas}
\begin{itemize}
  \item Anatomy: $\dot{\vb x}=\vb f(\vb x,t,\vb u;\boldsymbol\beta)$, $\vb y=\vb g(\vb x,t)$; state,
        dynamics (vector field), time, control, parameters, measurements.
  \item Linear systems: $\vb x(t)=Te^{\Lambda t}T^{-1}\vb x(0)=\sum_j b_je^{\lambda_jt}\boldsymbol\xi_j$;
        decoupled in $\vb z=T^{-1}\vb x$; discrete-time eigenvalues $\mu_j=e^{\lambda_j\Delta t}$.
  \item Primary challenges: nonlinearity, unknown $\vb f$. Secondary: dimension, multiscale, chaos,
        latent variables, noise, uncertainty. All recast as optimization.
  \item Uses: prediction (also of ensembles), design, control, interpretability and generalizability.
  \item Koopman program: find $\boldsymbol\varphi(\vb x)$ with $\tfrac{\mathrm d}{\mathrm dt}\boldsymbol\varphi=K\boldsymbol\varphi$.
\end{itemize}
\end{keyformulas}

% ------------------------------------------------------------------
\section*{Exercises}

Standalone exercises follow the section they practise; connected or integrative ones, like this one, are at the end.

\begin{exercise}[Anatomy of a system of your choice]\label{ex:ov-anatomy}
Choose a system: (a) an SIR epidemic in a city; (b) a wind turbine in a turbulent wind; (c) the
temperature of a room heated by a thermostat. For each one:
\begin{enumerate}[(i)]
  \item define a state $\vb x$ and estimate its dimension;
  \item list the controls $\vb u$ and the parameters $\boldsymbol\beta$, and say which parameters
        could become controls;
  \item say what can realistically be measured, $\vb y=\vb g(\vb x,t)$;
  \item say which challenges of \cref{sec:ov-challenges} apply, and which use of
        \cref{sec:ov-uses} you would care most about.
\end{enumerate}
\end{exercise}

\end{paracol}
